\documentclass[10pt,a4paper]{amsart}
\usepackage[margin=19mm]{geometry}
\usepackage{amsmath,amssymb,mathtools}
\usepackage[scr=rsfs,cal=euler]{mathalfa}
\usepackage{xfrac}
\usepackage[unicode,hidelinks,bookmarks=false]{hyperref}
\newcommand{\ie}{i.e.\ }
\newcommand{\cf}{cf.\ }
\newcommand{\resp}{respectively\ }
\newcommand{\Subsection}{\subsection}
\providecommand{\nobreakdash}{\nobreak\hbox{-}}
\setlength{\emergencystretch}{2em}
\multlinegap0pt
\usepackage{bm}
\usepackage{bbm}
\usepackage{dsfont}
\usepackage{xspace}
\usepackage{cancel}
\usepackage{mathtools}
\usepackage{amsthm}
\makeatletter
\@ifundefined{theorem}{\newtheorem{theorem}{Theorem}}{}
\@ifundefined{proposition}{\newtheorem{proposition}[theorem]{Proposition}}{}
\makeatother
\DeclareMathOperator{\Gal}{\mathrm {Gal}}
\DeclareMathOperator{\Hom}{\mathrm{Hom}}
\DeclareMathOperator{\Specm}{\mathrm{Specm}}
\title[Rational structures on quivers and a generalization of Gelfa]{Rational structures on quivers and a generalization of Gelfand's equivalence}
\author{Fabian Januszewski}
\date{Source: arXiv:2506.23251v1 (CC BY 4.0)}
\begin{document}
\maketitle

\noindent\emph{Source and licence.} Rational structures on quivers and a generalization of Gelfand's equivalence. Version arXiv:2506.23251v1 (CC BY 4.0), distributed under the Creative Commons Attribution 4.0 International licence, as recorded in the arXiv licence metadata for this exact version and shown on the abstract page https://arxiv.org/abs/2506.23251v1. Full rights notice: licenses/arxiv-2506.23251-CC-BY-4.0.txt.\par\smallskip

\noindent\textit{Editorial scope.} Counted unit: the Theorem whose source label is \texttt{thm:KrationalquiveretaleKspeciesantiequivalence} in the article (source0.tex lines 733--742, source label \texttt{thm:KrationalquiveretaleKspeciesantiequivalence}), delivered together with the article's own complete original proof of it (source0.tex lines 744--835, one \texttt{proof} environment of 92 lines). No mathematics is written, repaired or re-derived: statement and proof are the article's own text line for line. Required context retained uncounted: prop:Viorbits (lines 678--684). Every definition, display and label of those lines is retained unchanged. Omitted from this package and not redistributed: 1--677, 685--732, 743--743, 836--2622. Nothing inside a retained excerpt is omitted or altered: every retained body is delivered exactly as the article's lines are written, so no omission receipt of any kind accompanies this package. Citations: the retained excerpts carry no citation command at all, so no bibliography entry of the article is transcribed and no reference is redistributed. Reference adaptation: none was needed and none was made; every reference inside the delivered unit resolves inside it, so no unresolved reference or citation is produced. Printed slips kept literal: no printed word, symbol, hypothesis or number of the counted statement or of its proof was altered. Layout and class: the article's own class is \texttt{amsart} with packages including amsmath, amssymb, amsfonts, amscd, amsthm, mathrsfs, tikz, xy, mathalfa, natbib, geometry, hyperref; the wrapper is the lane's pinned \texttt{amsart} editorial wrapper at 10pt on A4, which restates the theorem-like environments the retained text needs and adds the article's own macro definitions that the retained text needs (\texttt{\textbackslash Gal}, \texttt{\textbackslash Hom}, \texttt{\textbackslash Specm}), with extra wrapper packages (bm, bbm, dsfont, xspace, cancel, mathtools). The delivered numbering is the unit's own. Assets: no retained span contains a figure environment, a graphics-inclusion command, a PDF-page inclusion, a table or a TikZ picture; no image file is copied into this package, the package declares no asset dependency and no image file is redistributed with it. Licence: version arXiv:2506.23251v1 of this article is distributed under the Creative Commons Attribution 4.0 International licence, as recorded in the arXiv licence metadata for this exact version and shown on the abstract page https://arxiv.org/abs/2506.23251v1. A full rights notice is delivered at licenses/arxiv-2506.23251-CC-BY-4.0.txt. Characters: every retained excerpt is pure ASCII, so the delivered engine, XeLaTeX with its default Latin Modern text font, reports no missing character.
  \begin{proposition}\label{prop:Viorbits}
  The fibers of the canonical map
  \[
  \pi\colon V\to I,\;\sigma\mapsto i,\text{ if }\sigma\in\Hom_K(L_i,L),
  \]
  are in canonical bijection with the $G$-orbits in $V$.
  \end{proposition}

  \begin{theorem}\label{thm:KrationalquiveretaleKspeciesantiequivalence}
    The maps
    \[
    S(-)\colon \{\text{\rm $K$-rational quivers}\}\to\{S\;\text{\rm \'etale $K$-species}\},
    \]
    \[
    \Gamma_K(-)\colon \{\text{\rm \'etale $K$-species}\}\to\{\text{\rm $K$-rational quivers}\},
    \]
    extend to quasi-inverse contravariant functors of the underlying categories, inducing anti-equivalences of the underlying full subcategories of split objects.
  \end{theorem}

  \begin{proof}
    The proof consists of two main parts. First, we show that for any $K$-rational quiver $\Gamma_K$, the canonical map $\Gamma_K \to \Gamma_K(S(\Gamma_K))$ is an isomorphism. Second, we show that for any \'etale $K$-species $S$, the canonical map $S \to S(\Gamma_K(S))$ is an isomorphism. Throughout the proof, we let $G = \Gal(L/K)$ with $L$ sufficiently big such that all actions are defined as actions of $\Gal(L/K)$ and all occuring finite \'etale algebras split over $L$. Alternatively, we may choose $L=K^{\rm sep}$ at the expense of working with an infinite Galois extension throughout.

    \paragraph{Step 1: $\Gamma_K(S(\Gamma_K)) \cong \Gamma_K$.}\ 

    \smallskip
    Let $\Gamma_K = (V,E,s,t)$ be a $K$-rational quiver with its $G$-action. We apply the functor $S(-)$ to obtain the \'etale $K$-species $S(\Gamma_K) = (L_i, {}_iM_j)_{i,j \in I}$. Recall that $I=G\backslash V$ is the set of $G$-orbits of vertices. For each orbit $i\in I$, we fix a representative $v_i \in i$, so $L_i=L(v_i)=L^{G_{v_i}}$. For any two orbits $i,j \in I$, the bimodule ${}_iM_j$ is defined as
    \[
      {}_iM_j = \bigoplus\limits_{\varepsilon \in G\backslash E_{ij}} L(e_\varepsilon),
    \]
    where $E_{ij}=\{e\in E \mid s(e)\in i, t(e)\in j\}$, and for each orbit $\varepsilon=Ge \in G\backslash E_{ij}$, we choose a representative $e_\varepsilon \in \varepsilon$, so $L(e_\varepsilon) = L^{G_{e_\varepsilon}}$.

    Now we apply the functor $\Gamma_K(-)$ to $S(\Gamma_K)$ to obtain a new $K$-rational quiver $\Gamma'_K = (V',E',s',t')$. We must construct a $G$-equivariant isomorphism $(\Phi_V, \Phi_E): \Gamma_K \to \Gamma'_K$.

    \smallskip
    \noindent{\bf Natural bijection on the vertices:} By definition,
    \[
    V' = \bigsqcup\limits_{i\in I} \Hom_K(L_i, L).
    \]
    For a given vertex $v \in V$, let $i=Gv$ be its orbit. We have $v=\tau v_i$ for some $\tau \in G$. We define a map $\Phi_V\colon V \to V'$ by sending $v$ to the $K$-algebra homomorphism $\phi_v \in \Hom_K(L_i,L)$ given by $\phi_v(\alpha) = \tau \alpha$ for all $\alpha\in L_i = L(v_i)$. This map is well-defined, since if $v=\tau'v_i$ for another $\tau'\in G$, then $\tau'=\tau \eta$ for some $\eta\in G_{v_i}$, and for $\alpha\in L(v_i)$, $\tau'(\alpha) = \tau \eta \alpha = \tau \alpha$. The map $\Phi_V$ is a bijection of $G$-sets. The inverse maps an embedding $\phi \in \Hom_K(L(v_i),L)$ to $\tau v_i$, where $\tau\in G$ is any element such that $\phi(\alpha)=\tau \alpha$ for all $\alpha\in L(v_i)$. Such a $\tau\in G$ exists and is unique up to left multiplication by an element of $G_{v_i}$. The $G$-equivariance is readily checked: for $\sigma\in G$, $\Phi_V(\sigma v)=\Phi_V(\sigma \tau v_i)$ is the map $\alpha \mapsto \sigma \tau \alpha$, which is precisely $\sigma \Phi_V(v)$.

    \smallskip
    \noindent{\bf Natural bijection on the edges:} By definition,
    \[
    E' = \bigsqcup\limits_{i,j\in I}\bigsqcup\limits_{\varepsilon'\in\Specm({}_iM_j)}\Hom_{K}(L_{\varepsilon'},L).
    \]
    The set of maximal ideals $\Specm({}_iM_j)$ of
    \[
      {}_iM_j = \bigoplus\limits_{\varepsilon \in G\backslash E_{ij}} L(e_\varepsilon)
    \]
    is in natural bijection with the set of edge orbits $G\backslash E_{ij}$. For an orbit $\varepsilon \in G\backslash E_{ij}$, the corresponding residue field is $L_{\varepsilon'} = L(e_\varepsilon)$. Thus, we can write
    \[
    E' = \bigsqcup\limits_{i,j\in I}\bigsqcup_{\varepsilon \in G\backslash E_{ij}} \Hom_{K}(L(e_\varepsilon),L).
    \]
    This is the disjoint union over all edge orbits in $G\backslash E$.

    For an edge $e\in E$, let $\varepsilon=Ge$ be its orbit. We have $e=\sigma e_\varepsilon$ for some $\sigma\in G$. We define $\Phi_E\colon E \to E'$ by sending $e$ to the $K$-algebra homomorphism $\phi_e \in \Hom_{K}(L(e_\varepsilon),L)$ given by $\phi_e(\alpha)=\sigma \alpha$. As with vertices, this map is well-defined and provides a $G$-equivariant bijection $E \to E'$.

    \smallskip
    \noindent{\bf Compatibility with source and target maps:} We must check that $\Phi_V(s(e)) = s'(\Phi_E(e))$ and $\Phi_V(t(e)) = t'(\Phi_E(e))$. Let $e=\sigma e_\varepsilon \in E$ for $\sigma\in G$, with $s(e_\varepsilon)\in i=Gv_i$ and $t(e_\varepsilon)\in j=Gv_j$. We may choose $v_i=s(e_\varepsilon)$ and $v_j=t(e_\varepsilon)$.
    
    The source map $s'$ for $\Gamma'_K$ is defined as $s'(\phi_e) = \phi_e \circ \tau_\varepsilon$, where $\tau_\varepsilon\colon L_i \to L(e_\varepsilon)$ is the structure map, which in this context is the inclusion $L(s(e_\varepsilon))\hookrightarrow L(e_\varepsilon)$.
So, for $\alpha\in L_i = L(s(e_\varepsilon))$, we have $s'(\phi_e)(\alpha) = \phi_e(\alpha) = \sigma \alpha$.

    On the other hand, $s(e) = s(\sigma e_\varepsilon)=\sigma s(e_\varepsilon)$. Thus $\Phi_V(s(e))$ is the map $\alpha \mapsto \sigma \alpha$. Hence $\Phi_V(s(e)) = s'(\Phi_E(e))$. The argument for $t$ is identical. This concludes the proof that $\Gamma_K(S(\Gamma_K))\cong \Gamma_K$.

    \paragraph{Step 2: $S(\Gamma_K(S)) \cong S$.}\ 

    \smallskip
Let $S=(L_i, {}_iM_j)_{i,j\in I}$ be an \'etale $K$-species. We construct the $K$-rational quiver $\Gamma_K(S)=(V,E,s,t)$. Then we apply $S(-)$ to obtain $S' = S(\Gamma_K(S))=(L'_{i'},{}_{i'}M'_{j'})_{i',j'\in I'}$. We need to construct an isomorphism of \'etale $K$-species from $S$ to $S'$.

    \smallskip
    \noindent{\bf Natural bijection on the index sets:} The index set of $S'$ is $I' = G\backslash V$. By construction,
    \[
    V=\bigsqcup\limits_{i\in I}\Hom_K(L_i,L).
    \]
    Each set $\Hom_K(L_i,L)$ is a single $G$-orbit in $V$. Thus, there is a natural bijection $\iota\colon I \to I'$ where $\iota(i)$ is the orbit $\Hom_K(L_i,L)$ (cf. Propositition \ref{prop:Viorbits}).

    \smallskip
    \noindent{\bf Natural isomorphisms of the fields:} For $i'=\iota(i)\in I'$, the field $L'_{i'}$ is $L^{G_{v'}}$ for a chosen representative $v' \in i'$. Let us choose $v'\colon L_i \to L$ to be an embedding. The stabilizer $G_{v'}$ consists of all $\sigma\in G$ such that $\sigma\circ v' = v'$. By Galois theory, the fixed field $L^{G_{v'}}$ is exactly the image $v'(L_i)$, which is $K$-isomorphic to $L_i$ via $v'$. Thus, we have $K$-algebra isomorphisms $\iota_i\colon L_i \to L'_{i'}$ for all $i$.

    \smallskip
    \noindent{\bf Natural isomorphisms of the bimodules (bialgebras):} For $i'=\iota(i), j'=\iota(j) \in I'$, the bimodule is ${}_{i'}M'_{j'} = \bigoplus_{\varepsilon'' \in G\backslash E_{i'j'}} L(e_{\varepsilon''})$.
    
    The set of edges $E_{i'j'}$ consists of all $\tau\in E$ with $s(\tau)\in i'$ and $t(\tau)\in j'$. By construction of
    \[
    E = \bigsqcup\limits_{k,l\in I}\bigsqcup\limits_{\varepsilon\in\Specm({}_kM_l)}\Hom_K(L_\varepsilon,L),
    \]
    an edge $\tau\in E$ has source in $\Hom_K(L_k,L)$ and target in $\Hom_K(L_l,L)$. Thus $E_{i'j'}$ consists of edges constructed from the bimodules ${}_iM_j$, i.e.,
    \[
    E_{i'j'} = \bigsqcup\limits_{\varepsilon\in\Specm({}_iM_j)}\Hom_K(L_\varepsilon,L).
    \]
    The $G$-orbits $\varepsilon''$ in $E_{i'j'}$ are precisely the sets $\Hom_K(L_\varepsilon,L)$ for $\varepsilon \in\Specm({}_iM_j)$.
    
    For such an orbit $\varepsilon''$, we pick a representative $e_{\varepsilon''}$, which is an embedding $\tau\colon L_\varepsilon \to L$. The field $L(e_{\varepsilon''}) = L^{G_{\tau}}$ is isomorphic to $L_\varepsilon$.
    
    The $L'_{i'},L'_{j'}$-bimodule structure on $L(e_{\varepsilon''}) \cong L_\varepsilon$ is induced by the maps $s(e_{\varepsilon''})$ and $t(e_{\varepsilon''})$. These are $s(\tau) = \tau\circ\tau_\varepsilon \in i'$ and $t(\tau) = \tau\circ\sigma_\varepsilon \in j'$, where $\tau_\varepsilon: L_i\to L_\varepsilon$ and $\sigma_\varepsilon: L_j\to L_\varepsilon$ are the algebra structure maps. This structure corresponds exactly to the original $L_i,L_j$-bimodule structure on $L_\varepsilon$. Therefore,
    \[
      {}_{i'}M'_{j'} = \bigoplus_{\varepsilon \in \Specm({}_iM_j)} L_\varepsilon.
    \]
    Since ${}_iM_j$ is a finite \'etale $K$-algebra, it decomposes into a product of fields
    \[
      {}_iM_j = \prod\limits_{\varepsilon \in \Specm({}_iM_j)} L_\varepsilon.
    \]
    For a finite index set, the direct sum and direct product of bimodules coincide. Thus we have an isomorphism of $L_i,L_j$-bimodules (and $K$-bialgebras) $\iota_{ij}\colon {}_iM_j \to {}_{i'}M'_{j'}$.
    
    This collection of maps $(\iota, (\iota_i)_i, (\iota_{ij})_{i,j})$ defines an isomorphism of \'etale $K$-species $S \to S(\Gamma_K(S))$.
    
    \paragraph{Functoriality and split objects:} The constructions are natural. A morphism of quivers induces a morphism of species in the opposite direction and vice versa by the well known anti-equivalence between $G$-sets and finite commutative \'etale $K$-algebras split over $L$.
    
    If a quiver $\Gamma_K$ is split over $K$, its $G$-action is trivial. Then $G_v=G$ for all $v\in V$ and $G_e=G$ for all $e\in E$. Consequently, $L_i=L^G=K$ and $L(e_\varepsilon)=K$ for all $i$ and $\varepsilon$. The resulting species $S(\Gamma_K)$ is a split \'etale $K$-species. Conversely, starting from a split species $S$ (with all $L_i=K$ and ${}_iM_j$ a product of copies of $K$), the resulting quiver $\Gamma_K(S)$ will have trivial $G$-action on its vertices and edges, making it a split $K$-rational quiver. The anti-equivalence thus restricts to an anti-equivalence between the full subcategories of split objects.
  \end{proof}

\end{document}
